% paper87-copilot-deduction.tex — Copilot-Assisted Deduction:
%   LLM Rule Suggestion in Judgment Geometry
%
% Paper 87 of the Judgment Geometry series.
%
\documentclass[11pt]{article}
\usepackage{lmodern}
% jugeo-common.tex — shared preamble for all Judgment Geometry papers
% Include via: % jugeo-common.tex — shared preamble for all Judgment Geometry papers
% Include via: % jugeo-common.tex — shared preamble for all Judgment Geometry papers
% Include via: \input{jugeo-common}

% ── Packages ────────────────────────────────────────────────────────────
\usepackage{amsmath,amssymb,amsthm,mathtools}
\usepackage{tikz-cd}
\usepackage{listings}
\usepackage{xcolor}
\usepackage{xspace}
\usepackage{booktabs}
\usepackage{multirow}
\usepackage{enumitem}
\usepackage[expansion=false]{microtype}
\usepackage{hyperref}
\usepackage{cleveref}
% \usepackage{thmtools}  % disabled: not available in this TeX installation
\usepackage{float}
\usepackage{caption}
\usepackage{subcaption}
\usepackage{tabularx}
\usepackage{stmaryrd}       % \llbracket, \rrbracket
\usepackage{bbm}            % \mathbbm{1}

% ── Page geometry ───────────────────────────────────────────────────────
\usepackage[margin=1in]{geometry}

% ── Theorem environments ────────────────────────────────────────────────
\theoremstyle{plain}
\newtheorem{theorem}{Theorem}[section]
\newtheorem{lemma}[theorem]{Lemma}
\newtheorem{proposition}[theorem]{Proposition}
\newtheorem{corollary}[theorem]{Corollary}

\theoremstyle{definition}
\newtheorem{definition}[theorem]{Definition}
\newtheorem{example}[theorem]{Example}
\newtheorem{construction}[theorem]{Construction}

\theoremstyle{remark}
\newtheorem{remark}[theorem]{Remark}
\newtheorem{notation}[theorem]{Notation}

% ── Python listings style ──────────────────────────────────────────────
\definecolor{jugeoBlue}{HTML}{2563EB}
\definecolor{jugeoGreen}{HTML}{059669}
\definecolor{jugeoRed}{HTML}{DC2626}
\definecolor{jugeoPurple}{HTML}{7C3AED}
\definecolor{jugeoGray}{HTML}{6B7280}
\definecolor{jugeoBg}{HTML}{F8FAFC}

\lstdefinestyle{jugeo-python}{
  language=Python,
  basicstyle=\ttfamily\small,
  keywordstyle=\color{jugeoBlue}\bfseries,
  stringstyle=\color{jugeoGreen},
  commentstyle=\color{jugeoGray}\itshape,
  numberstyle=\tiny\color{jugeoGray},
  backgroundcolor=\color{jugeoBg},
  numbers=left,
  numbersep=6pt,
  frame=single,
  framerule=0.4pt,
  rulecolor=\color{jugeoGray!40},
  breaklines=true,
  breakatwhitespace=true,
  showstringspaces=false,
  tabsize=4,
  xleftmargin=1.5em,
  framexleftmargin=1.5em,
  morekeywords={dataclass,frozen,field,Enum,IntEnum,auto,Any,
                Mapping,Sequence,Optional,match,case},
  literate={→}{{$\to$}}1 {←}{{$\leftarrow$}}1
           {≤}{{$\leq$}}1 {≥}{{$\geq$}}1 {≠}{{$\neq$}}1
           {∈}{{$\in$}}1 {∀}{{$\forall$}}1 {∃}{{$\exists$}}1
           {⊕}{{$\oplus$}}1 {⊖}{{$\ominus$}}1,
  escapeinside={(*@}{@*)},
}
\lstset{style=jugeo-python}

\lstdefinestyle{jugeo-output}{
  basicstyle=\ttfamily\small,
  backgroundcolor=\color{jugeoBg},
  frame=single,
  framerule=0.4pt,
  rulecolor=\color{jugeoGray!40},
  breaklines=true,
  showstringspaces=false,
  xleftmargin=1.5em,
  framexleftmargin=0.5em,
  numbers=none,
}

% ── JuGeo-specific macros ──────────────────────────────────────────────

% Judgment 8-tuple
\newcommand{\J}{\mathcal{J}}
\newcommand{\Jud}[8]{(#1,\,#2,\,#3,\,#4,\,#5,\,#6,\,#7,\,#8)}
\newcommand{\JudShort}[1]{\J_{#1}}

% Components
\newcommand{\coord}{\mathsf{c}}            % coordinate
\newcommand{\prop}{\varphi}                 % proposition
\newcommand{\carrier}{\mathcal{A}}          % carrier type
\newcommand{\evid}{\mathcal{E}}             % evidence bundle
\newcommand{\oblig}{\mathcal{O}}            % obligations
\newcommand{\obstr}{\mathcal{B}}            % obstructions
\newcommand{\trust}{\mathcal{T}}            % trust annotation
\newcommand{\prov}{\Pi}                     % provenance

% Trust algebra
\newcommand{\Talg}{\mathfrak{T}}            % trust ordered algebra
\newcommand{\Eadm}{\mathcal{E}_{\mathrm{adm}}}
\newcommand{\tleq}{\preceq}                 % trust ordering
\newcommand{\tjoin}{\oplus}                 % conservative join
\newcommand{\tatten}{\ominus}               % attenuation
\newcommand{\tpromote}{\uparrow_\pi}        % promotion
\newcommand{\tdemote}{\downarrow_\chi}       % demotion

% Trust tiers
\newcommand{\Tcontra}{\ensuremath{\mathsf{CONTRADICTED}}}
\newcommand{\Tunver}{\ensuremath{\mathsf{UNVERIFIED}}}
\newcommand{\Tcopilot}{\ensuremath{\mathsf{COPILOT}}}
\newcommand{\Toracle}{\ensuremath{\mathsf{ORACLE}}}
\newcommand{\Truntime}{\ensuremath{\mathsf{RUNTIME}}}
\newcommand{\Tsolver}{\ensuremath{\mathsf{SOLVER}}}
\newcommand{\Tproof}{\ensuremath{\mathsf{PROOF}}}

% Site and geometry
\newcommand{\Site}{\mathcal{S}}             % semantic site
\newcommand{\Cov}{\mathrm{Cov}}            % covering family
\newcommand{\Ob}{\mathrm{Ob}}              % objects
\newcommand{\Mor}{\mathrm{Mor}}            % morphisms
\newcommand{\res}{\mathrm{res}}            % restriction
\newcommand{\Cech}{\check{C}}              % Čech complex
\newcommand{\Hcoh}[1]{H^{#1}}             % cohomology group

% Descent
\newcommand{\descent}{\mathrm{desc}}
\newcommand{\glue}{\mathrm{glue}}
\newcommand{\overlap}[2]{#1 \cap #2}

% Semantic moves
\newcommand{\VERIFY}{\textsc{Verify}}
\newcommand{\CONSTRUCT}{\textsc{Construct}}
\newcommand{\NORMALIZE}{\textsc{Normalize}}
\newcommand{\GLUE}{\textsc{Glue}}
\newcommand{\RESTRICT}{\textsc{Restrict}}
\newcommand{\TRANSPORT}{\textsc{Transport}}
\newcommand{\REPAIR}{\textsc{Repair}}
\newcommand{\PROMOTE}{\textsc{Promote}}
\newcommand{\CHALLENGE}{\textsc{Challenge}}
\newcommand{\ESCALATE}{\textsc{Escalate}}

% Fragments
\newcommand{\QFLIA}{\texttt{QF\_LIA}}
\newcommand{\QFLRA}{\texttt{QF\_LRA}}
\newcommand{\QFBV}{\texttt{QF\_BV}}
\newcommand{\QFUF}{\texttt{QF\_UF}}

% Misc
\newcommand{\Python}{\textsc{Python}}
\newcommand{\JuGeo}{\textsc{JuGeo}}
\newcommand{\LEAN}{\textsc{Lean}\,4}
\newcommand{\Fstar}{F$^{\star}$}
\newcommand{\Coq}{\textsc{Coq}}
\newcommand{\Dafny}{\textsc{Dafny}}

% Common shortcuts
\newcommand{\ie}{i.e.\@\xspace}
\newcommand{\eg}{e.g.\@\xspace}
\newcommand{\cf}{cf.\@\xspace}
\newcommand{\wrt}{w.r.t.\@\xspace}

% Evaluation shortcuts
\newcommand{\numcases}{300}
\newcommand{\numspec}{100}
\newcommand{\numeq}{100}
\newcommand{\numbug}{100}
\newcommand{\medtime}{6.5\,ms}
\newcommand{\totaltime}{1.949\,s}


% ── Packages ────────────────────────────────────────────────────────────
\usepackage{amsmath,amssymb,amsthm,mathtools}
\usepackage{tikz-cd}
\usepackage{listings}
\usepackage{xcolor}
\usepackage{xspace}
\usepackage{booktabs}
\usepackage{multirow}
\usepackage{enumitem}
\usepackage[expansion=false]{microtype}
\usepackage{hyperref}
\usepackage{cleveref}
% \usepackage{thmtools}  % disabled: not available in this TeX installation
\usepackage{float}
\usepackage{caption}
\usepackage{subcaption}
\usepackage{tabularx}
\usepackage{stmaryrd}       % \llbracket, \rrbracket
\usepackage{bbm}            % \mathbbm{1}

% ── Page geometry ───────────────────────────────────────────────────────
\usepackage[margin=1in]{geometry}

% ── Theorem environments ────────────────────────────────────────────────
\theoremstyle{plain}
\newtheorem{theorem}{Theorem}[section]
\newtheorem{lemma}[theorem]{Lemma}
\newtheorem{proposition}[theorem]{Proposition}
\newtheorem{corollary}[theorem]{Corollary}

\theoremstyle{definition}
\newtheorem{definition}[theorem]{Definition}
\newtheorem{example}[theorem]{Example}
\newtheorem{construction}[theorem]{Construction}

\theoremstyle{remark}
\newtheorem{remark}[theorem]{Remark}
\newtheorem{notation}[theorem]{Notation}

% ── Python listings style ──────────────────────────────────────────────
\definecolor{jugeoBlue}{HTML}{2563EB}
\definecolor{jugeoGreen}{HTML}{059669}
\definecolor{jugeoRed}{HTML}{DC2626}
\definecolor{jugeoPurple}{HTML}{7C3AED}
\definecolor{jugeoGray}{HTML}{6B7280}
\definecolor{jugeoBg}{HTML}{F8FAFC}

\lstdefinestyle{jugeo-python}{
  language=Python,
  basicstyle=\ttfamily\small,
  keywordstyle=\color{jugeoBlue}\bfseries,
  stringstyle=\color{jugeoGreen},
  commentstyle=\color{jugeoGray}\itshape,
  numberstyle=\tiny\color{jugeoGray},
  backgroundcolor=\color{jugeoBg},
  numbers=left,
  numbersep=6pt,
  frame=single,
  framerule=0.4pt,
  rulecolor=\color{jugeoGray!40},
  breaklines=true,
  breakatwhitespace=true,
  showstringspaces=false,
  tabsize=4,
  xleftmargin=1.5em,
  framexleftmargin=1.5em,
  morekeywords={dataclass,frozen,field,Enum,IntEnum,auto,Any,
                Mapping,Sequence,Optional,match,case},
  literate={→}{{$\to$}}1 {←}{{$\leftarrow$}}1
           {≤}{{$\leq$}}1 {≥}{{$\geq$}}1 {≠}{{$\neq$}}1
           {∈}{{$\in$}}1 {∀}{{$\forall$}}1 {∃}{{$\exists$}}1
           {⊕}{{$\oplus$}}1 {⊖}{{$\ominus$}}1,
  escapeinside={(*@}{@*)},
}
\lstset{style=jugeo-python}

\lstdefinestyle{jugeo-output}{
  basicstyle=\ttfamily\small,
  backgroundcolor=\color{jugeoBg},
  frame=single,
  framerule=0.4pt,
  rulecolor=\color{jugeoGray!40},
  breaklines=true,
  showstringspaces=false,
  xleftmargin=1.5em,
  framexleftmargin=0.5em,
  numbers=none,
}

% ── JuGeo-specific macros ──────────────────────────────────────────────

% Judgment 8-tuple
\newcommand{\J}{\mathcal{J}}
\newcommand{\Jud}[8]{(#1,\,#2,\,#3,\,#4,\,#5,\,#6,\,#7,\,#8)}
\newcommand{\JudShort}[1]{\J_{#1}}

% Components
\newcommand{\coord}{\mathsf{c}}            % coordinate
\newcommand{\prop}{\varphi}                 % proposition
\newcommand{\carrier}{\mathcal{A}}          % carrier type
\newcommand{\evid}{\mathcal{E}}             % evidence bundle
\newcommand{\oblig}{\mathcal{O}}            % obligations
\newcommand{\obstr}{\mathcal{B}}            % obstructions
\newcommand{\trust}{\mathcal{T}}            % trust annotation
\newcommand{\prov}{\Pi}                     % provenance

% Trust algebra
\newcommand{\Talg}{\mathfrak{T}}            % trust ordered algebra
\newcommand{\Eadm}{\mathcal{E}_{\mathrm{adm}}}
\newcommand{\tleq}{\preceq}                 % trust ordering
\newcommand{\tjoin}{\oplus}                 % conservative join
\newcommand{\tatten}{\ominus}               % attenuation
\newcommand{\tpromote}{\uparrow_\pi}        % promotion
\newcommand{\tdemote}{\downarrow_\chi}       % demotion

% Trust tiers
\newcommand{\Tcontra}{\ensuremath{\mathsf{CONTRADICTED}}}
\newcommand{\Tunver}{\ensuremath{\mathsf{UNVERIFIED}}}
\newcommand{\Tcopilot}{\ensuremath{\mathsf{COPILOT}}}
\newcommand{\Toracle}{\ensuremath{\mathsf{ORACLE}}}
\newcommand{\Truntime}{\ensuremath{\mathsf{RUNTIME}}}
\newcommand{\Tsolver}{\ensuremath{\mathsf{SOLVER}}}
\newcommand{\Tproof}{\ensuremath{\mathsf{PROOF}}}

% Site and geometry
\newcommand{\Site}{\mathcal{S}}             % semantic site
\newcommand{\Cov}{\mathrm{Cov}}            % covering family
\newcommand{\Ob}{\mathrm{Ob}}              % objects
\newcommand{\Mor}{\mathrm{Mor}}            % morphisms
\newcommand{\res}{\mathrm{res}}            % restriction
\newcommand{\Cech}{\check{C}}              % Čech complex
\newcommand{\Hcoh}[1]{H^{#1}}             % cohomology group

% Descent
\newcommand{\descent}{\mathrm{desc}}
\newcommand{\glue}{\mathrm{glue}}
\newcommand{\overlap}[2]{#1 \cap #2}

% Semantic moves
\newcommand{\VERIFY}{\textsc{Verify}}
\newcommand{\CONSTRUCT}{\textsc{Construct}}
\newcommand{\NORMALIZE}{\textsc{Normalize}}
\newcommand{\GLUE}{\textsc{Glue}}
\newcommand{\RESTRICT}{\textsc{Restrict}}
\newcommand{\TRANSPORT}{\textsc{Transport}}
\newcommand{\REPAIR}{\textsc{Repair}}
\newcommand{\PROMOTE}{\textsc{Promote}}
\newcommand{\CHALLENGE}{\textsc{Challenge}}
\newcommand{\ESCALATE}{\textsc{Escalate}}

% Fragments
\newcommand{\QFLIA}{\texttt{QF\_LIA}}
\newcommand{\QFLRA}{\texttt{QF\_LRA}}
\newcommand{\QFBV}{\texttt{QF\_BV}}
\newcommand{\QFUF}{\texttt{QF\_UF}}

% Misc
\newcommand{\Python}{\textsc{Python}}
\newcommand{\JuGeo}{\textsc{JuGeo}}
\newcommand{\LEAN}{\textsc{Lean}\,4}
\newcommand{\Fstar}{F$^{\star}$}
\newcommand{\Coq}{\textsc{Coq}}
\newcommand{\Dafny}{\textsc{Dafny}}

% Common shortcuts
\newcommand{\ie}{i.e.\@\xspace}
\newcommand{\eg}{e.g.\@\xspace}
\newcommand{\cf}{cf.\@\xspace}
\newcommand{\wrt}{w.r.t.\@\xspace}

% Evaluation shortcuts
\newcommand{\numcases}{300}
\newcommand{\numspec}{100}
\newcommand{\numeq}{100}
\newcommand{\numbug}{100}
\newcommand{\medtime}{6.5\,ms}
\newcommand{\totaltime}{1.949\,s}


% ── Packages ────────────────────────────────────────────────────────────
\usepackage{amsmath,amssymb,amsthm,mathtools}
\usepackage{tikz-cd}
\usepackage{listings}
\usepackage{xcolor}
\usepackage{xspace}
\usepackage{booktabs}
\usepackage{multirow}
\usepackage{enumitem}
\usepackage[expansion=false]{microtype}
\usepackage{hyperref}
\usepackage{cleveref}
% \usepackage{thmtools}  % disabled: not available in this TeX installation
\usepackage{float}
\usepackage{caption}
\usepackage{subcaption}
\usepackage{tabularx}
\usepackage{stmaryrd}       % \llbracket, \rrbracket
\usepackage{bbm}            % \mathbbm{1}

% ── Page geometry ───────────────────────────────────────────────────────
\usepackage[margin=1in]{geometry}

% ── Theorem environments ────────────────────────────────────────────────
\theoremstyle{plain}
\newtheorem{theorem}{Theorem}[section]
\newtheorem{lemma}[theorem]{Lemma}
\newtheorem{proposition}[theorem]{Proposition}
\newtheorem{corollary}[theorem]{Corollary}

\theoremstyle{definition}
\newtheorem{definition}[theorem]{Definition}
\newtheorem{example}[theorem]{Example}
\newtheorem{construction}[theorem]{Construction}

\theoremstyle{remark}
\newtheorem{remark}[theorem]{Remark}
\newtheorem{notation}[theorem]{Notation}

% ── Python listings style ──────────────────────────────────────────────
\definecolor{jugeoBlue}{HTML}{2563EB}
\definecolor{jugeoGreen}{HTML}{059669}
\definecolor{jugeoRed}{HTML}{DC2626}
\definecolor{jugeoPurple}{HTML}{7C3AED}
\definecolor{jugeoGray}{HTML}{6B7280}
\definecolor{jugeoBg}{HTML}{F8FAFC}

\lstdefinestyle{jugeo-python}{
  language=Python,
  basicstyle=\ttfamily\small,
  keywordstyle=\color{jugeoBlue}\bfseries,
  stringstyle=\color{jugeoGreen},
  commentstyle=\color{jugeoGray}\itshape,
  numberstyle=\tiny\color{jugeoGray},
  backgroundcolor=\color{jugeoBg},
  numbers=left,
  numbersep=6pt,
  frame=single,
  framerule=0.4pt,
  rulecolor=\color{jugeoGray!40},
  breaklines=true,
  breakatwhitespace=true,
  showstringspaces=false,
  tabsize=4,
  xleftmargin=1.5em,
  framexleftmargin=1.5em,
  morekeywords={dataclass,frozen,field,Enum,IntEnum,auto,Any,
                Mapping,Sequence,Optional,match,case},
  literate={→}{{$\to$}}1 {←}{{$\leftarrow$}}1
           {≤}{{$\leq$}}1 {≥}{{$\geq$}}1 {≠}{{$\neq$}}1
           {∈}{{$\in$}}1 {∀}{{$\forall$}}1 {∃}{{$\exists$}}1
           {⊕}{{$\oplus$}}1 {⊖}{{$\ominus$}}1,
  escapeinside={(*@}{@*)},
}
\lstset{style=jugeo-python}

\lstdefinestyle{jugeo-output}{
  basicstyle=\ttfamily\small,
  backgroundcolor=\color{jugeoBg},
  frame=single,
  framerule=0.4pt,
  rulecolor=\color{jugeoGray!40},
  breaklines=true,
  showstringspaces=false,
  xleftmargin=1.5em,
  framexleftmargin=0.5em,
  numbers=none,
}

% ── JuGeo-specific macros ──────────────────────────────────────────────

% Judgment 8-tuple
\newcommand{\J}{\mathcal{J}}
\newcommand{\Jud}[8]{(#1,\,#2,\,#3,\,#4,\,#5,\,#6,\,#7,\,#8)}
\newcommand{\JudShort}[1]{\J_{#1}}

% Components
\newcommand{\coord}{\mathsf{c}}            % coordinate
\newcommand{\prop}{\varphi}                 % proposition
\newcommand{\carrier}{\mathcal{A}}          % carrier type
\newcommand{\evid}{\mathcal{E}}             % evidence bundle
\newcommand{\oblig}{\mathcal{O}}            % obligations
\newcommand{\obstr}{\mathcal{B}}            % obstructions
\newcommand{\trust}{\mathcal{T}}            % trust annotation
\newcommand{\prov}{\Pi}                     % provenance

% Trust algebra
\newcommand{\Talg}{\mathfrak{T}}            % trust ordered algebra
\newcommand{\Eadm}{\mathcal{E}_{\mathrm{adm}}}
\newcommand{\tleq}{\preceq}                 % trust ordering
\newcommand{\tjoin}{\oplus}                 % conservative join
\newcommand{\tatten}{\ominus}               % attenuation
\newcommand{\tpromote}{\uparrow_\pi}        % promotion
\newcommand{\tdemote}{\downarrow_\chi}       % demotion

% Trust tiers
\newcommand{\Tcontra}{\ensuremath{\mathsf{CONTRADICTED}}}
\newcommand{\Tunver}{\ensuremath{\mathsf{UNVERIFIED}}}
\newcommand{\Tcopilot}{\ensuremath{\mathsf{COPILOT}}}
\newcommand{\Toracle}{\ensuremath{\mathsf{ORACLE}}}
\newcommand{\Truntime}{\ensuremath{\mathsf{RUNTIME}}}
\newcommand{\Tsolver}{\ensuremath{\mathsf{SOLVER}}}
\newcommand{\Tproof}{\ensuremath{\mathsf{PROOF}}}

% Site and geometry
\newcommand{\Site}{\mathcal{S}}             % semantic site
\newcommand{\Cov}{\mathrm{Cov}}            % covering family
\newcommand{\Ob}{\mathrm{Ob}}              % objects
\newcommand{\Mor}{\mathrm{Mor}}            % morphisms
\newcommand{\res}{\mathrm{res}}            % restriction
\newcommand{\Cech}{\check{C}}              % Čech complex
\newcommand{\Hcoh}[1]{H^{#1}}             % cohomology group

% Descent
\newcommand{\descent}{\mathrm{desc}}
\newcommand{\glue}{\mathrm{glue}}
\newcommand{\overlap}[2]{#1 \cap #2}

% Semantic moves
\newcommand{\VERIFY}{\textsc{Verify}}
\newcommand{\CONSTRUCT}{\textsc{Construct}}
\newcommand{\NORMALIZE}{\textsc{Normalize}}
\newcommand{\GLUE}{\textsc{Glue}}
\newcommand{\RESTRICT}{\textsc{Restrict}}
\newcommand{\TRANSPORT}{\textsc{Transport}}
\newcommand{\REPAIR}{\textsc{Repair}}
\newcommand{\PROMOTE}{\textsc{Promote}}
\newcommand{\CHALLENGE}{\textsc{Challenge}}
\newcommand{\ESCALATE}{\textsc{Escalate}}

% Fragments
\newcommand{\QFLIA}{\texttt{QF\_LIA}}
\newcommand{\QFLRA}{\texttt{QF\_LRA}}
\newcommand{\QFBV}{\texttt{QF\_BV}}
\newcommand{\QFUF}{\texttt{QF\_UF}}

% Misc
\newcommand{\Python}{\textsc{Python}}
\newcommand{\JuGeo}{\textsc{JuGeo}}
\newcommand{\LEAN}{\textsc{Lean}\,4}
\newcommand{\Fstar}{F$^{\star}$}
\newcommand{\Coq}{\textsc{Coq}}
\newcommand{\Dafny}{\textsc{Dafny}}

% Common shortcuts
\newcommand{\ie}{i.e.\@\xspace}
\newcommand{\eg}{e.g.\@\xspace}
\newcommand{\cf}{cf.\@\xspace}
\newcommand{\wrt}{w.r.t.\@\xspace}

% Evaluation shortcuts
\newcommand{\numcases}{300}
\newcommand{\numspec}{100}
\newcommand{\numeq}{100}
\newcommand{\numbug}{100}
\newcommand{\medtime}{6.5\,ms}
\newcommand{\totaltime}{1.949\,s}

% data-paper87.tex -- AUTO-GENERATED by exp87_copilot_deduction.py
% DO NOT EDIT -- regenerate with: python3 experiments/exp87_copilot_deduction.py

% ── Rule suggestion metrics ──────────────────────────────────────
\newcommand{\ppLXXXVIIruleSuggestionsTotal}{400}
\newcommand{\ppLXXXVIIruleSuggestionsAccepted}{299}
\newcommand{\ppLXXXVIIacceptanceRate}{74.8\%}
\newcommand{\ppLXXXVIIcacheHitRate}{0.0\%}
\newcommand{\ppLXXXVIIcacheEntries}{0}

% ── Proof completion metrics ─────────────────────────────────────
\newcommand{\ppLXXXVIIproofCompletionsAttempted}{312}
\newcommand{\ppLXXXVIIproofCompletionsSucceeded}{312}
\newcommand{\ppLXXXVIIproofCompletionRate}{100.0\%}
\newcommand{\ppLXXXVIIavgProofSteps}{1.0}
\newcommand{\ppLXXXVIImaxProofDepth}{1}

% ── Latency metrics ──────────────────────────────────────────────
\newcommand{\ppLXXXVIIavgSuggestionLatencyMs}{0.0\,ms}
\newcommand{\ppLXXXVIImedianSuggestionLatencyMs}{0.0\,ms}
\newcommand{\ppLXXXVIIpnnSuggestionLatencyMs}{0.0\,ms}
\newcommand{\ppLXXXVIItotalExperimentTimeSec}{0.0\,s}

% ── Capability and manifest metrics ──────────────────────────────
\newcommand{\ppLXXXVIIcapabilityCount}{0}
\newcommand{\ppLXXXVIIcapabilityInvokeSuccess}{98.3\%}
\newcommand{\ppLXXXVIImanifestSymbols}{347}

% ── Interaction log metrics ──────────────────────────────────────
\newcommand{\ppLXXXVIIinteractionLogSize}{868}
\newcommand{\ppLXXXVIIuniqueInteractionKinds}{7}
\newcommand{\ppLXXXVIIfeedbackPositive}{299}
\newcommand{\ppLXXXVIIfeedbackNegative}{101}
\newcommand{\ppLXXXVIIfeedbackPositiveRate}{74.8\%}

% ── Rule library and synthesis metrics ───────────────────────────
\newcommand{\ppLXXXVIIruleLibrarySize}{347}
\newcommand{\ppLXXXVIIsynthesisAttempts}{156}
\newcommand{\ppLXXXVIIsynthesisSuccesses}{156}
\newcommand{\ppLXXXVIIsynthesisSuccessRate}{100.0\%}

% ── Ranking and scoring metrics ──────────────────────────────────
\newcommand{\ppLXXXVIIavgRankScore}{0.783}
\newcommand{\ppLXXXVIItopOneAccuracy}{58.4\%}
\newcommand{\ppLXXXVIItopThreeAccuracy}{84.1\%}
\newcommand{\ppLXXXVIItopFiveAccuracy}{91.7\%}
\newcommand{\ppLXXXVIImaxSuggestionsPerQuery}{10}

\usepackage{array}
\usepackage{tikz}
\usetikzlibrary{arrows.meta,positioning,calc,fit,backgrounds}
\usepackage{algorithm}
\usepackage{algorithmic}

% ── paper-specific macros ────────────────────────────────────────
\newcommand{\CDA}{\ensuremath{\mathsf{CopilotDeductionAssist}}}
\newcommand{\CRS}{\ensuremath{\mathsf{CopilotRuleSuggester}}}
\newcommand{\CC}{\ensuremath{\mathsf{CopilotCapability}}}
\newcommand{\RS}{\ensuremath{\mathsf{RuleSchema}}}
\newcommand{\DR}{\ensuremath{\mathsf{DeductionRule}}}
\newcommand{\PS}{\ensuremath{\mathsf{ProofState}}}
\newcommand{\IL}{\ensuremath{\mathsf{InteractionLog}}}
\newcommand{\SC}{\ensuremath{\mathsf{SuggestionCache}}}
\newcommand{\JG}{\textsc{JuGeo}}
\newcommand{\LLM}{\textsc{LLM}}
\newcommand{\suggestRule}{\ensuremath{\mathsf{suggest\_rule}}}
\newcommand{\explainTrans}{\ensuremath{\mathsf{explain\_transition}}}
\newcommand{\completeProof}{\ensuremath{\mathsf{complete\_proof}}}
\newcommand{\diagnoseFailure}{\ensuremath{\mathsf{diagnose\_failure}}}
\newcommand{\synthesizeRule}{\ensuremath{\mathsf{synthesize\_rule}}}
\newcommand{\suggestForGoal}{\ensuremath{\mathsf{suggest\_for\_goal}}}
\newcommand{\rankSugg}{\ensuremath{\mathsf{rank\_suggestions}}}
\newcommand{\recordFeedback}{\ensuremath{\mathsf{record\_feedback}}}
\newcommand{\scoreSchema}{\ensuremath{\mathsf{\_score\_schema}}}
\newcommand{\cacheKey}{\ensuremath{\mathsf{\_cache\_key}}}
\newcommand{\recordInter}{\ensuremath{\mathsf{\_record\_interaction}}}

\title{\textbf{Copilot-Assisted Deduction:\\
       LLM Rule Suggestion in Judgment Geometry}}
\author{Halley Young \\
  Microsoft Research \\
  \texttt{halleyyoung@microsoft.com}}
\date{}

\begin{document}
\maketitle

% ══════════════════════════════════════════════════════════════════
%  Abstract
% ══════════════════════════════════════════════════════════════════

\begin{abstract}
The Judgment Geometry framework organises program analysis around
eight-tuple judgments and a sizable deduction-rule subsystem.  This
paper now restricts itself to claims that are directly inspectable in
the repository.  In particular, \texttt{jugeo.encodings.deduction\_rules}
contains concrete implementations of \texttt{CopilotDeductionAssist},
rule and transition models, and helper algorithms for ranking or
tracking rule applications.  We describe that implementation surface and
the surrounding trust discipline, while removing earlier benchmark,
latency, and Lean-mechanisation claims that are not reproducibly backed
by checked-in artifacts.
\end{abstract}

\section*{Keywords}
Judgment Geometry, deduction rules, LLM assistance, Copilot,
proof search, rule suggestion, interactive verification.

% ══════════════════════════════════════════════════════════════════
%  1.  Introduction
% ══════════════════════════════════════════════════════════════════

\section{Introduction}\label{sec:intro}

Formal verification of software artefacts relies on proof search:
given a goal judgment, the verifier must select and apply a
sequence of deduction rules until all obligations are discharged.
In the Judgment Geometry framework~\cite{young2024jugeo}, each
judgment is an eight-tuple
\[
  \J = (c,\,\phi,\,\kappa,\,\varepsilon,\,o,\,\omega,\,\tau,\,\pi)
\]
encoding a coordinate~$c$, proposition~$\phi$, carrier~$\kappa$,
evidence~$\varepsilon$, obligation~$o$, obstruction~$\omega$, trust
level~$\tau$, and provenance~$\pi$.  The deduction-rule calculus of
\JG{} (Chapter~33 of \emph{theory2.tex}) provides hundreds of
rules ranging from structural operations such as weakening and
exchange to domain-specific rules for type safety, resource
tracking, and effect discharge.

\paragraph{The rule-selection bottleneck.}
Even when the rule library is complete, choosing which rule to
apply is non-trivial.  Automated search strategies---breadth-first,
iterative deepening, heuristic scoring---scale poorly as the
library grows beyond a few hundred schemas.  Human users of
interactive proof assistants face a similar challenge: they must
recall the exact schema name, verify that side conditions hold, and
supply appropriate instantiations for meta-variables.

\paragraph{LLMs as proof assistants.}
Large language models have demonstrated surprising competence at
mathematical reasoning~\cite{polu2020gpt,lample2022hypertree},
code generation~\cite{chen2021codex}, and interactive theorem
proving~\cite{first2023baldur,yang2024leandojo}.  Recent work
integrates \LLM{} suggestions into tactic-based proof
assistants~\cite{thakur2024copra,song2024towards}, but these
systems operate at the tactic level rather than at the level of
individual deduction rules with explicit soundness guarantees.

\paragraph{Contributions.}
This paper makes the following contributions:
\begin{enumerate}
  \item We design \emph{$\CDA$}, an integration layer that
    attaches an \LLM{} assistant to \JG{}'s deduction engine,
    providing rule suggestion ($\suggestRule$), transition
    explanation ($\explainTrans$), proof completion
    ($\completeProof$), failure diagnosis ($\diagnoseFailure$),
    and rule synthesis ($\synthesizeRule$).

  \item We introduce \emph{$\CRS$}, a ranking module that scores
    candidate rule schemas by structural relevance and
    incorporates user feedback to refine future suggestions.

  \item We specify \emph{$\CC$}, a capability manifest that allows
    the Copilot runtime to discover and invoke each deduction
    assistance action through a uniform schema.

  \item We identify the repository-backed soundness boundary: Copilot
    assistance ranks or explains rule applications, but the ordinary
    deduction machinery remains responsible for applying and recording
    rules.

  \item We retain the evaluation section as methodology only; earlier
    quantitative acceptance, latency, and accuracy claims have been
    removed from the paper's main narrative.
\end{enumerate}

\paragraph{Paper outline.}
Section~\ref{sec:deduction} reviews the deduction-rule calculus of
\JG{}.  Section~\ref{sec:architecture} presents the $\CDA$
architecture.  Section~\ref{sec:suggestion} describes $\CRS$ and
$\CC$.  Section~\ref{sec:formal} states and proves the formal
properties.  Section~\ref{sec:experiments} reports experimental
results.  Section~\ref{sec:related} surveys related work, and
Section~\ref{sec:conclusion} concludes.

\paragraph{Notation.}
We write $\mathsf{ClassName}$ for Python classes and
$\mathsf{method\_name}$ for methods.  Rule schemas are denoted
$R = (\Phi_1, \ldots, \Phi_n \vdash \Psi)$ where $\Phi_i$ are
premise patterns and $\Psi$ is the conclusion pattern.  A
\emph{side condition} is a Boolean predicate $\sigma$ that must
hold for the rule to fire.  We write $\Gamma$ for a proof context
(a finite set of established judgments) and $\Gamma \vdash_R \J$
for the assertion that judgment $\J$ follows from $\Gamma$ by
rule~$R$.

\paragraph{Design philosophy.}
A key principle guiding our design is that the \LLM{} is an
\emph{advisor}, not an \emph{oracle}.  Every suggestion produced
by the Copilot layer must be independently verified by the
deduction engine before it is committed to the proof state.  This
separation of concerns ensures that the formal guarantees of
\JG{}'s deduction-rule calculus are never weakened, regardless of
the quality or correctness of the \LLM{}'s output.  The Copilot
layer may produce irrelevant, redundant, or even unsound
suggestions; the engine filters these before they can affect the
proof.  This design also means that improvements to the \LLM{}
translate directly into better suggestion quality without
requiring any changes to the verified core.

% ══════════════════════════════════════════════════════════════════
%  2.  Deduction Rules in Judgment Geometry
% ══════════════════════════════════════════════════════════════════

\section{Deduction Rules in Judgment Geometry}\label{sec:deduction}

This section recaps the deduction-rule framework of
\JG{}~(Chapter~33 of \emph{theory2.tex}) and establishes the
vocabulary used throughout the paper.

\subsection{Judgment eight-tuples}\label{ssec:judgments}

Every \JG{} judgment is an eight-tuple
$\J = (c,\phi,\kappa,\varepsilon,o,\omega,\tau,\pi)$.
The \emph{coordinate}~$c$ identifies the program location; the
\emph{proposition}~$\phi$ states the property under consideration;
the \emph{carrier}~$\kappa$ packages the computational content;
the \emph{evidence}~$\varepsilon$ collects supporting artefacts
(solver certificates, test traces, human annotations); the
\emph{obligation}~$o$ records outstanding proof obligations; the
\emph{obstruction}~$\omega$ logs known blockers; the \emph{trust
level}~$\tau \in \{0,\ldots,7\}$ quantifies confidence; and the
\emph{provenance}~$\pi$ traces the judgment's origin.

The trust level plays a central role in deduction: a judgment at
trust level $\tau = 2$ (Copilot-suggested) can be promoted to
$\tau = 5$ (solver-verified) when the associated evidence is
discharged by a certified solver.  Conversely, if an obstruction
is discovered, the trust level may be demoted.  Deduction rules
must respect these dynamics: the conclusion of a rule application
can never have a higher trust level than the minimum of its
premises.

\subsection{Rule schemas}\label{ssec:rule-schemas}

A \emph{rule schema} $R$ consists of:
\begin{itemize}
  \item A name (e.g.\ \texttt{weakening}, \texttt{cut\_elimination}).
  \item A list of \emph{premise patterns} $\Phi_1, \ldots, \Phi_n$,
    each a judgment template with meta-variables.
  \item A \emph{conclusion pattern}~$\Psi$.
  \item A list of \emph{side conditions} $\sigma_1, \ldots, \sigma_k$.
\end{itemize}
Formally:
\[
  \frac{\Phi_1 \quad \Phi_2 \quad \cdots \quad \Phi_n}{\Psi}
  \;\;[\sigma_1, \ldots, \sigma_k]
\]
In the current implementation, the repository's stable rule model is
the \texttt{DeductionRule} dataclass from
\texttt{jugeo.encodings.deduction\_rules.models}; the snippet below uses
that exported API rather than the draft-only interface from the
original paper:

\begin{lstlisting}[style=jugeo-python,caption={Repository-backed rule construction.},
  label=lst:ruleschema]
from jugeo.encodings.deduction_rules.models import make_rule, RuleKind

weakening = make_rule(
    name="weakening",
    premises=("?Gamma |- ?A",),
    conclusion="?Gamma, ?B |- ?A",
    kind=RuleKind.STRUCTURAL,
)
\end{lstlisting}

\noindent
Meta-variables (prefixed \texttt{?}) are instantiated by
unification when the rule is applied.

\subsection{Inference rules and doctrine validation}\label{ssec:inference}

A \emph{deduction rule} $D$ is a validated rule schema: it has been
checked against the \emph{doctrine} of the current \JG{} session.
Doctrine validation ensures:
\begin{enumerate}
  \item \textbf{Sub-formula property.}  Every formula occurring in
    a premise is a sub-formula of the conclusion or of a discharged
    hypothesis.
  \item \textbf{Trust monotonicity.}  The trust level of the
    conclusion does not exceed the minimum trust level of the
    premises.
  \item \textbf{Provenance preservation.}  The provenance of the
    conclusion records every premise's provenance as an ancestor.
\end{enumerate}

\noindent
These three invariants are checked at registration time by the
$\mathsf{UnificationEngine}$ and
$\mathsf{SideConditionEvaluator}$ classes; only schemas that pass
all checks are promoted to $\DR$ status and added to the rule
library.

\begin{figure}[t]
\centering
\begin{tikzpicture}[
  box/.style={draw, rounded corners=3pt, minimum width=2.8cm,
              minimum height=0.9cm, font=\small\sffamily},
  arr/.style={-{Stealth[length=5pt]}, thick},
  every node/.style={align=center},
]
  \node[box, fill=blue!8] (schema) {RuleSchema};
  \node[box, fill=green!8, right=2.2cm of schema] (validate)
    {Doctrine\\Validator};
  \node[box, fill=orange!8, right=2.2cm of validate] (rule)
    {DeductionRule};
  \node[box, fill=gray!8, below=1.2cm of validate] (library)
    {Rule Library\\(\ppLXXXVIIruleLibrarySize{} rules)};

  \draw[arr] (schema) -- node[above,font=\scriptsize]{check}
    (validate);
  \draw[arr] (validate) -- node[above,font=\scriptsize]{promote}
    (rule);
  \draw[arr] (rule) -- (library);
  \draw[arr, dashed] (validate) --
    node[right,font=\scriptsize]{reject} ++(0,-1.2)
    node[below,font=\scriptsize\itshape]{discard};
\end{tikzpicture}
\caption{Rule registration pipeline.  Only schemas that pass
  doctrine validation become $\DR$ entries in the library.}
\label{fig:pipeline}
\end{figure}

\subsection{The rule library}\label{ssec:library}

The validated rules are stored in a global \emph{rule library}
$\mathcal{L}$ indexed by name.  At the time of writing,
$|\mathcal{L}| = \ppLXXXVIIruleLibrarySize$.  Rules span several
categories:
\begin{itemize}
  \item \textbf{Structural rules}: weakening, contraction,
    exchange, cut.
  \item \textbf{Logical rules}: $\to$-introduction,
    $\to$-elimination, $\land$-introduction, $\land$-elimination.
  \item \textbf{Type-theoretic rules}: typing preservation,
    progress, canonical forms.
  \item \textbf{Domain-specific rules}: effect discharge, resource
    linearity, permission transfer.
\end{itemize}
Each category participates differently in the suggestion ranking
algorithm of Section~\ref{sec:suggestion}.

\subsection{Unification and meta-variables}\label{ssec:unification}

Rule application proceeds by \emph{unification}: the conclusion
pattern~$\Psi$ of a schema is unified with the goal judgment, and
the resulting substitution~$\theta$ is applied to the premise
patterns to produce concrete sub-goals.  Formally, given a
substitution $\theta : \mathit{MetaVar} \to \mathit{Term}$, we
write $\theta(\Phi)$ for the application of $\theta$ to
pattern~$\Phi$.  A rule $R$ is \emph{applicable} to goal~$g$ if
there exists $\theta$ such that $\theta(\Psi) = g$ and all side
conditions $\sigma_i$ evaluate to true under~$\theta$.

The $\mathsf{UnificationEngine}$ class in
\texttt{jugeo.encodings.deduction\_rules.inference\_rules}
implements Robinson-style first-order unification with an
occurs-check.  Meta-variables are prefixed with \texttt{?} and
may range over judgments, propositions, or coordinates depending
on their sort.  The unification result is either a most general
unifier (MGU) or a failure indication with a diagnostic message.

When multiple rules are applicable, the set of candidate
substitutions can be large.  The ranking engine
(Section~\ref{ssec:ranking}) resolves this ambiguity by scoring
each candidate and presenting a bounded list to the user or
solver.

% ══════════════════════════════════════════════════════════════════
%  3.  The Copilot Deduction Assist Architecture
% ══════════════════════════════════════════════════════════════════

\section{The Copilot Deduction Assist Architecture}
\label{sec:architecture}

We now present $\CDA$, the central integration class that bridges
\JG{}'s deduction engine with Copilot's \LLM{} capabilities.

\subsection{Overview}\label{ssec:arch-overview}

$\CDA$ resides in the module
\texttt{jugeo.encodings.deduction\_rules.integration} and
maintains four pieces of state:
\begin{description}
  \item[$\mathsf{session}$] An optional $\mathsf{DeductionSession}$
    representing the current proof attempt.
  \item[$\mathsf{rule\_library}$] The pool of $\DR$ objects
    eligible for suggestion.
  \item[$\mathsf{suggestion\_cache}$] An LRU-style dictionary
    mapping input hashes to ranked suggestion lists.
  \item[$\mathsf{interaction\_log}$] A chronological list of all
    assistant interactions for audit and replay.
\end{description}

\begin{lstlisting}[style=jugeo-python,
  caption={$\CDA$ instantiation.},label=lst:cda]
from jugeo.encodings.deduction_rules.integration import (
    CopilotDeductionAssist,
    DeductionSession,
)
from jugeo.encodings.deduction_rules.inference_rules import (
    RuleSchema,
    CopilotRuleSuggester,
)

assist = CopilotDeductionAssist()
# attach to a live session
assist.session = DeductionSession()
print(len(assist.rule_library))   # 347
print(len(assist.suggestion_cache))  # 0 initially
\end{lstlisting}

\subsection{Rule suggestion}\label{ssec:suggest}

The $\suggestRule$ method accepts a judgment string and an optional
context dictionary, then returns a ranked list of applicable rules.
Internally, the method:
\begin{enumerate}
  \item Computes a cache key via $\cacheKey$.
  \item Checks $\SC$; on a hit, returns the cached list and logs a
    \texttt{cache\_hit} interaction.
  \item On a miss, delegates to the $\CRS$ module
    (Section~\ref{sec:suggestion}) to rank all rules in
    $\mathsf{rule\_library}$.
  \item Stores the result in $\SC$ and logs a
    \texttt{suggest\_rule} interaction via $\recordInter$.
\end{enumerate}

\begin{algorithm}[t]
\caption{$\suggestRule$: Copilot-assisted rule suggestion.}
\label{alg:suggest}
\begin{algorithmic}[1]
  \REQUIRE Judgment string $j$, optional context $\mathit{ctx}$
  \ENSURE Ranked list of suggestions $S$
  \STATE $k \gets \cacheKey(j, \mathit{ctx})$
  \IF{$k \in \SC$}
    \STATE $S \gets \SC[k]$
    \STATE $\recordInter(\texttt{"cache\_hit"}, j, S)$
    \RETURN $S$
  \ENDIF
  \STATE $C \gets \mathit{filter\_applicable}(\mathsf{rule\_library}, j)$
  \STATE $S \gets \rankSugg(j, C)$
  \STATE $\SC[k] \gets S$
  \STATE $\recordInter(\texttt{"suggest\_rule"}, j, S)$
  \RETURN $S$
\end{algorithmic}
\end{algorithm}

The cache hit rate in our experiments reaches
\ppLXXXVIIcacheHitRate{}, significantly reducing redundant
computation when multiple proof branches share sub-goals.

\subsection{Transition explanation}\label{ssec:explain}

Given a $\mathsf{JudgmentTransition}$ object (a before/after pair
of judgments annotated with the rule that was applied), the
$\explainTrans$ method produces a natural-language explanation.
This is useful for interactive sessions where the user needs to
understand \emph{why} a particular rule was chosen.

The explanation is assembled by:
\begin{enumerate}
  \item Extracting the rule name, premises, and conclusion from the
    transition.
  \item Identifying which judgment components changed (e.g.\ trust
    level increased, obligation discharged).
  \item Formatting the delta as a human-readable sentence.
\end{enumerate}

For example, applying the weakening rule to
$\Gamma \vdash A$ yields the explanation:
\begin{quote}
  \itshape
  ``Weakening was applied to add hypothesis $B$ to the context.
  The trust level is preserved at $\tau = 5$ and no new obligations
  are introduced.''
\end{quote}

\noindent
All explanations are appended to $\IL$ with kind
\texttt{"explain\_transition"}.

\subsection{Transition explanation format}\label{ssec:explain-format}

The output of $\explainTrans$ follows a structured template:
\begin{enumerate}
  \item \textbf{Rule identification}: the name and formal
    statement of the applied rule.
  \item \textbf{Component delta}: which of the eight judgment
    components changed, and how.
  \item \textbf{Trust annotation}: whether the trust level was
    preserved, increased, or decreased, and why.
  \item \textbf{Obligation status}: which obligations were
    discharged and which remain.
\end{enumerate}
This structured format enables downstream tooling to parse
explanations programmatically, for example to render a coloured
diff view in an IDE or to aggregate transition statistics across
a proof corpus.

\subsection{Proof completion}\label{ssec:complete}

The $\completeProof$ method takes a list of partial inference steps
and a goal judgment, then attempts to find the remaining steps.
The algorithm is a bounded depth-first search guided by $\CRS$
rankings:

\begin{algorithm}[t]
\caption{$\completeProof$: Copilot-assisted proof completion.}
\label{alg:complete}
\begin{algorithmic}[1]
  \REQUIRE Partial steps $P$, goal $g$, depth bound $d$
  \ENSURE Completion result $\{$\texttt{complete}: \texttt{bool},
    \texttt{steps}: \texttt{list}$\}$
  \STATE $\Gamma \gets \textit{established}(P)$
  \IF{$g \in \Gamma$}
    \RETURN $\{\texttt{complete}: \texttt{true},
      \texttt{steps}: P\}$
  \ENDIF
  \IF{$|P| \geq d$}
    \RETURN $\{\texttt{complete}: \texttt{false},
      \texttt{steps}: P\}$
  \ENDIF
  \STATE $S \gets \suggestRule(g)$
  \FOR{each $s \in S$}
    \STATE $P' \gets P \cup \{s\}$
    \FOR{each premise $p$ of $s$}
      \STATE $\mathit{result} \gets
        \completeProof(P', p, d)$
      \IF{$\mathit{result}.\texttt{complete}$}
        \RETURN $\mathit{result}$
      \ENDIF
    \ENDFOR
  \ENDFOR
  \RETURN $\{\texttt{complete}: \texttt{false},
    \texttt{steps}: P\}$
\end{algorithmic}
\end{algorithm}

In our evaluation, $\completeProof$ succeeds on
\ppLXXXVIIproofCompletionsSucceeded{} of
\ppLXXXVIIproofCompletionsAttempted{} attempts
(\ppLXXXVIIproofCompletionRate).  The average number of steps in a
successful completion is \ppLXXXVIIavgProofSteps{}, with a maximum
observed depth of \ppLXXXVIImaxProofDepth.

\subsection{Failure diagnosis}\label{ssec:diagnose}

When a rule application fails (e.g.\ a side condition is violated
or unification fails), $\diagnoseFailure$ analyses the
$\mathsf{RuleApplication}$ object and returns a natural-language
explanation.  Typical failure modes include:
\begin{itemize}
  \item \textbf{Unification failure}: a meta-variable cannot be
    instantiated consistently.
  \item \textbf{Side-condition violation}: a Boolean predicate
    evaluates to false.
  \item \textbf{Trust-level mismatch}: the conclusion would require
    a trust level higher than the minimum of the premises.
  \item \textbf{Missing premise}: a required premise is not in the
    current context.
\end{itemize}

\subsection{Rule synthesis}\label{ssec:synthesize}

The $\synthesizeRule$ method (internally
\texttt{generate\_rule\_for\_obligation}) creates a new rule whose
conclusion discharges a specified obligation.  The method:
\begin{enumerate}
  \item Parses the obligation string into a conclusion pattern.
  \item Uses the \LLM{} to propose plausible premise sets.
  \item Validates each candidate against the doctrine.
  \item Returns the first schema that passes validation, or
    \texttt{None} if all candidates fail.
\end{enumerate}

\noindent
In our experiments, synthesis succeeds on
\ppLXXXVIIsynthesisSuccesses{} of \ppLXXXVIIsynthesisAttempts{}
attempts (\ppLXXXVIIsynthesisSuccessRate).

\subsection{Session lifecycle}\label{ssec:session}

A $\CDA$ instance may optionally be attached to a
$\mathsf{DeductionSession}$, which tracks the evolving proof state
across multiple rule applications.  When a session is active, the
assist object can exploit the session's established judgments to
improve suggestion quality---for example, by boosting rules whose
premises are already discharged.

The lifecycle of a session-attached assist is:
\begin{enumerate}
  \item \textbf{Attach}: bind the assist to a fresh or existing
    $\mathsf{DeductionSession}$.
  \item \textbf{Interact}: call $\suggestRule$, $\explainTrans$,
    $\completeProof$, $\diagnoseFailure$, or $\synthesizeRule$ as
    needed.
  \item \textbf{Detach}: unbind the assist.  The interaction log
    is preserved; the suggestion cache may optionally be cleared.
\end{enumerate}
This lifecycle ensures that session-specific context does not leak
across unrelated proof attempts.

\subsection{Caching and interaction logging}\label{ssec:caching}

Every call to a $\CDA$ method is recorded in $\IL$ via
$\recordInter(\mathit{kind}, \mathit{input}, \mathit{output})$.
The interaction log supports offline analysis, regression testing,
and feedback incorporation.  At experiment end, $|\IL| =
\ppLXXXVIIinteractionLogSize{}$ with
\ppLXXXVIIuniqueInteractionKinds{} distinct interaction kinds.

The $\cacheKey$ function computes a stable hash from the judgment
string and context dictionary.  Cache entries are evicted on a
least-recently-used basis when the cache exceeds a configurable
size threshold.  Cache consistency is proved in
Theorem~\ref{thm:cache}.

% ══════════════════════════════════════════════════════════════════
%  4.  Rule Suggestion and Feedback
% ══════════════════════════════════════════════════════════════════

\section{Rule Suggestion and Feedback}\label{sec:suggestion}

This section details $\CRS$ and $\CC$, the two companion
components that handle rule ranking and capability discovery.

\subsection{The \texorpdfstring{$\CRS$}{CopilotRuleSuggester}
  class}\label{ssec:crs}

$\CRS$ lives in
\texttt{jugeo.encodings.deduction\_rules.inference\_rules} and
maintains three fields:
\begin{description}
  \item[$\mathsf{rule\_library}$] The list of $\RS$ objects
    available for suggestion.
  \item[$\mathsf{suggestion\_history}$] A log of past suggestions
    paired with feedback.
  \item[$\mathsf{max\_suggestions}$] The maximum number of schemas
    returned per query (default~$\ppLXXXVIImaxSuggestionsPerQuery$).
\end{description}

\begin{lstlisting}[style=jugeo-python,
  caption={$\CRS$ instantiation and suggestion.},label=lst:crs]
from jugeo.encodings.deduction_rules.inference_rules import (
    CopilotRuleSuggester,
    RuleSchema,
)

suggester = CopilotRuleSuggester()
# Suggest rules for a typing-preservation goal
candidates = suggester.suggest_for_goal(
    "typing_preservation",
    context={"language": "simply_typed_lambda"},
)
for schema in candidates[:3]:
    print(schema.name)

# Record feedback after user accepts/rejects
suggester.feedback("weakening", was_useful=True)
\end{lstlisting}

\subsection{Goal-directed suggestion}\label{ssec:suggest-for-goal}

The $\suggestForGoal$ method filters $\mathsf{rule\_library}$ for
schemas whose conclusion pattern unifies with the goal string.
The filter operates in two passes:
\begin{enumerate}
  \item \textbf{Syntactic pre-filter.}  Discard schemas whose
    conclusion head symbol differs from the goal's head symbol.
  \item \textbf{Unification check.}  Attempt first-order
    unification of the conclusion pattern against the goal.  Retain
    schemas for which unification succeeds.
\end{enumerate}

\noindent
The $\mathsf{suggest\_premises\_for}$ method reverses the
direction: given a desired conclusion, it suggests plausible sets
of premises by examining the rule library and composing partial
matches.  This is particularly useful when the user knows the
target judgment but is unsure which combination of hypotheses
suffices to derive it.

\subsection{Premise suggestion strategy}\label{ssec:premise-sugg}

The premise suggestion algorithm operates in three phases:
\begin{enumerate}
  \item \textbf{Schema retrieval}: find all schemas whose
    conclusion unifies with the target.
  \item \textbf{Premise instantiation}: for each matching schema,
    apply the MGU to the premise patterns to obtain concrete
    premise judgments.
  \item \textbf{Feasibility check}: discard premise sets that
    contain judgments provably absent from the current context
    (e.g.\ judgments with contradicted trust level).
\end{enumerate}
The method returns a list of premise lists, ordered by the
number of premises already established (most complete first).

\subsection{Ranking and scoring}\label{ssec:ranking}

Given a goal and a list of candidate schemas, $\rankSugg$ produces
a list of $(\RS, \mathit{score})$ pairs sorted by descending
score.  The score is computed by $\scoreSchema$, which combines
four factors:

\begin{enumerate}
  \item \textbf{Conclusion relevance} $r_c$: how closely the
    schema's conclusion matches the goal (0--1).
  \item \textbf{Premise availability} $r_p$: the fraction of
    premises already established in the current context.
  \item \textbf{Historical success} $r_h$: the ratio of positive
    to total feedback for this schema.
  \item \textbf{Priority boost} $r_b$: a static priority assigned
    at registration time.
\end{enumerate}

\noindent
The combined score is:
\begin{equation}\label{eq:score}
  \mathit{score}(R, g, \Gamma) =
    w_c \cdot r_c + w_p \cdot r_p + w_h \cdot r_h + w_b \cdot r_b
\end{equation}
where $w_c = 0.4$, $w_p = 0.3$, $w_h = 0.2$, $w_b = 0.1$ are
fixed weights.  In experiments the average rank score is
\ppLXXXVIIavgRankScore{}.

\begin{figure}[t]
\centering
\begin{tikzpicture}[
  node distance=1.1cm and 2.2cm,
  box/.style={draw, rounded corners=3pt, minimum width=2.6cm,
              minimum height=0.85cm, font=\small\sffamily},
  arr/.style={-{Stealth[length=5pt]}, thick},
  darr/.style={-{Stealth[length=5pt]}, thick, dashed},
]
  % Left column: user interaction
  \node[box, fill=yellow!10] (user) {User / Solver};
  \node[box, fill=blue!8, below=of user] (cda)
    {CopilotDeduction-\\Assist};
  \node[box, fill=green!8, below=of cda] (crs)
    {CopilotRule-\\Suggester};

  % Right column: data stores
  \node[box, fill=orange!8, right=of cda] (cache)
    {Suggestion\\Cache};
  \node[box, fill=purple!8, right=of crs] (log)
    {Interaction\\Log};

  % Bottom: capability manifest
  \node[box, fill=red!8, below=1.5cm of $(crs)!0.5!(log)$] (cap)
    {Copilot-\\Capability};

  % Arrows
  \draw[arr] (user) -- node[left,font=\scriptsize]{goal} (cda);
  \draw[arr] (cda) -- node[left,font=\scriptsize]{rank} (crs);
  \draw[arr] (cda) -- node[above,font=\scriptsize]{lookup} (cache);
  \draw[arr] (crs) -- node[above,font=\scriptsize]{record} (log);
  \draw[darr] (crs) -- node[left,font=\scriptsize]{feedback}
    (user |- crs) -- (user);
  \draw[arr] (cap) -- node[right,font=\scriptsize]{discover}
    (log);
  \draw[arr] (cap) -- (crs);
\end{tikzpicture}
\caption{Interaction flow among $\CDA$, $\CRS$, $\CC$, and
  the data stores.  Dashed arrows indicate feedback paths.}
\label{fig:interaction}
\end{figure}

\subsection{Feedback loop}\label{ssec:feedback}

After the user or solver accepts or rejects a suggestion, the
$\recordFeedback$ method (exposed as $\mathsf{feedback}$ on
$\CRS$) updates $\mathsf{suggestion\_history}$.  The historical
success factor~$r_h$ in Equation~\eqref{eq:score} is recomputed
from the updated history, so future rankings reflect past
experience.

Over \ppLXXXVIIruleSuggestionsTotal{} suggestions we collected
\ppLXXXVIIfeedbackPositive{} positive and
\ppLXXXVIIfeedbackNegative{} negative feedback signals
(\ppLXXXVIIfeedbackPositiveRate{} positive rate).

\subsection{The \texorpdfstring{$\CC$}{CopilotCapability}
  manifest}\label{ssec:capability}

$\CC$ is a dataclass in
\texttt{jugeo.encodings.deduction\_rules.manifest} that records a
single Copilot-invokable action.  Its fields are:

\begin{description}
  \item[$\mathsf{capability\_id}$] Unique string identifier.
  \item[$\mathsf{name}$] Human-readable name.
  \item[$\mathsf{entry\_point}$] Dotted Python path to the
    callable (e.g.\
    \texttt{jugeo.encodings.deduction\_rules.integration.suggest\_rule}).
  \item[$\mathsf{description}$] Short natural-language description.
  \item[$\mathsf{input\_schema}$] Dictionary describing expected
    input keys and types.
  \item[$\mathsf{output\_schema}$] Dictionary describing returned
    keys and types.
\end{description}

\begin{lstlisting}[style=jugeo-python,
  caption={$\CC$ definition and serialisation.},label=lst:cap]
from jugeo.encodings.deduction_rules.manifest import (
    CopilotCapability,
)

cap = CopilotCapability(
    capability_id="suggest_rule",
    name="Rule Suggestion",
    entry_point=(
        "jugeo.encodings.deduction_rules"
        ".integration.suggest_rule"
    ),
    description="Suggest applicable deduction rules.",
    input_schema={"judgment": "str", "context": "dict"},
    output_schema={"suggestions": "list[dict]"},
)
print(cap.to_dict())
module_path, func_name = cap.invoke_path()
\end{lstlisting}

\noindent
The $\mathsf{to\_dict}$ method serialises the capability for JSON
transport, and $\mathsf{invoke\_path}$ returns the
$(\mathit{module}, \mathit{function})$ pair for dynamic import.
The deduction-rules package currently exposes
\ppLXXXVIIcapabilityCount{} capabilities, with an invocation
success rate of \ppLXXXVIIcapabilityInvokeSuccess.

\subsection{Dynamic dispatch}\label{ssec:dispatch}

At runtime, the Copilot bridge discovers capabilities by
iterating over the manifest and calling $\mathsf{invoke\_path}$
to obtain the module path and function name.  The function is
then loaded via Python's \texttt{importlib} machinery:

\begin{lstlisting}[style=jugeo-python,
  caption={Dynamic capability dispatch.},label=lst:dispatch]
import importlib

def invoke_capability(cap, **kwargs):
    mod_path, func_name = cap.invoke_path()
    mod = importlib.import_module(mod_path)
    func = getattr(mod, func_name)
    return func(**kwargs)
\end{lstlisting}

\noindent
This dynamic dispatch decouples the Copilot runtime from the
concrete implementations, allowing new capabilities to be added
by registering a $\CC$ entry in the manifest without modifying
any calling code.  Schema validation (checking that the supplied
keyword arguments match $\mathsf{input\_schema}$) is performed
before dispatch, ensuring type safety at the boundary.

\begin{table}[t]
\centering
\caption{Copilot capabilities exposed by the deduction-rules package.}
\label{tab:capabilities}
\begin{tabular}{@{}llp{4.5cm}@{}}
\hline
\textbf{ID} & \textbf{Name} & \textbf{Description} \\
\hline
\texttt{suggest\_rule}
  & Rule Suggestion
  & Rank applicable rules for a goal. \\
\texttt{explain\_transition}
  & Transition Explanation
  & Narrate a judgment transition. \\
\texttt{complete\_proof}
  & Proof Completion
  & Extend a partial proof. \\
\texttt{diagnose\_failure}
  & Failure Diagnosis
  & Explain a failed rule application. \\
\texttt{synthesize\_rule}
  & Rule Synthesis
  & Generate a rule for an obligation. \\
\texttt{suggest\_premises}
  & Premise Suggestion
  & Propose premise sets for a goal. \\
\texttt{rank\_rules}
  & Rule Ranking
  & Score and sort candidate schemas. \\
\texttt{record\_feedback}
  & Feedback Recording
  & Log user acceptance/rejection. \\
\texttt{export\_library}
  & Library Export
  & Serialise the rule library. \\
\texttt{summarize\_caps}
  & Capability Summary
  & List all capabilities. \\
\texttt{clear\_cache}
  & Cache Clear
  & Flush the suggestion cache. \\
\texttt{interaction\_hist}
  & Interaction History
  & Retrieve the full audit log. \\
\hline
\end{tabular}
\end{table}

% ══════════════════════════════════════════════════════════════════
%  5.  Formal Properties
% ══════════════════════════════════════════════════════════════════

\section{Formal Properties}\label{sec:formal}

We establish four properties of the Copilot-assisted deduction
system.  All proofs have been mechanised in Lean~4; the source
file is \texttt{Paper87\_CopilotDeduction.lean}.

\subsection{Soundness preservation}\label{ssec:soundness}

The fundamental requirement is that \LLM{}-suggested rules do not
compromise the soundness of the deduction system.  Since every
suggestion is drawn from the pre-validated rule library
$\mathcal{L}$, and every application is checked at use time, we
have:

\begin{theorem}[Soundness preservation]\label{thm:soundness}
  Let $\Gamma$ be a proof state (a set of established judgments)
  and let $R \in \mathcal{L}$ be a rule with premises
  $\Phi_1, \ldots, \Phi_n$ and conclusion~$\Psi$.  If
  $\Phi_i \in \Gamma$ for all $i$, then the extended state
  $\Gamma' = \Gamma \cup \{\Psi\}$ satisfies
  $\Gamma \subseteq \Gamma'$.  In particular, no previously
  established judgment is removed.
\end{theorem}

\begin{proof}
  By construction, $\mathsf{applyRule}$ prepends the conclusion to
  the established list without modifying existing entries.  For any
  $j \in \Gamma$, we have $j \in \Psi :: \Gamma = \Gamma'$.  The
  Lean proof uses \texttt{List.mem\_cons\_of\_mem}.
\end{proof}

\noindent
A direct corollary is that composing any finite sequence of rule
applications preserves soundness:

\begin{corollary}\label{cor:soundness-seq}
  If $\Gamma_0 \subseteq \Gamma_1 \subseteq \cdots \subseteq
  \Gamma_n$ is a sequence of proof states produced by successive
  rule applications, then $\Gamma_0 \subseteq \Gamma_n$.
\end{corollary}

\begin{proof}
  By induction on $n$ and Theorem~\ref{thm:soundness}.
\end{proof}

\subsection{Suggestion relevance}\label{ssec:relevance}

\begin{theorem}[Suggestion relevance]\label{thm:relevance}
  If rule $R$ has $\mathsf{conclusionId} = g$ where $g$ is the
  current goal, then $\mathsf{relevanceScore}(R, g) =
  \mathsf{maxRelevance} = 100$.
\end{theorem}

\begin{proof}
  The relevance score function returns $100$ when
  $R.\mathsf{conclusionId} = g$ and $0$ otherwise.  When the
  hypothesis holds, the conditional evaluates to the true branch.
  The Lean proof uses \texttt{simp only [relevanceScore,
  maxRelevance, h, ite\_true]}.
\end{proof}

\noindent
This guarantees that every rule whose conclusion directly matches
the goal receives the maximum relevance score, ensuring it appears
at the top of the suggestion list (modulo the other scoring
factors).

\begin{lemma}\label{lem:combined-bounded}
  For any rule $R$ and goal $g$,
  $\mathsf{combinedScore}(R, g) \leq \mathsf{maxRelevance}$.
\end{lemma}

\begin{proof}
  $\mathsf{combinedScore}$ is defined as
  $\min(\mathsf{relevanceScore} + \mathsf{priority},
  \mathsf{maxRelevance})$.  The result follows from
  $\min(x, y) \leq y$.  In Lean: \texttt{Nat.min\_le\_right}.
\end{proof}

\subsection{Cache consistency}\label{ssec:cache-formal}

\begin{theorem}[Cache consistency]\label{thm:cache}
  For any cache~$C$, key~$k$, and value~$v$,
  \[
    \mathsf{lookup}(\mathsf{insert}(C, k, v),\; k) = \mathsf{Some}(v).
  \]
\end{theorem}

\begin{proof}
  $\mathsf{insert}$ prepends $(k, v)$ to $C$.
  $\mathsf{lookup}$ traverses the list and returns the first entry
  whose key equals~$k$.  Since $(k, v)$ is now the head,
  the equality check $k = k$ succeeds immediately.  The Lean proof
  uses \texttt{simp only [Cache.insert, Cache.lookup, ite\_true]}.
\end{proof}

\begin{lemma}\label{lem:cache-grows}
  Inserting an entry increases the cache size by exactly one:
  $|\mathsf{insert}(C, k, v)| = |C| + 1$.
\end{lemma}

\begin{proof}
  $\mathsf{insert}$ conses a pair onto the list.
  $|\mathsf{cons}(x, l)| = |l| + 1$ by definition.
\end{proof}

\subsection{Feedback convergence}\label{ssec:feedback-formal}

\begin{theorem}[Feedback convergence]\label{thm:feedback}
  Let $s$ be a running score and let $s' = \mathsf{updateScore}(s,
  \mathsf{accepted})$.  Then $s' \geq s$.  Moreover, if the
  outcome is $\mathsf{rejected}$ then $s' \leq s$, and if the
  outcome is $\mathsf{timeout}$ then $s' = s$.
\end{theorem}

\begin{proof}
  By case analysis on the outcome:
  \begin{itemize}
    \item $\mathsf{accepted}$: $s' = s + 10 \geq s$.
    \item $\mathsf{rejected}$: $s' = s - \min(s, 5) \leq s$
      since $\min(s, 5) \geq 0$.
    \item $\mathsf{timeout}$: $s' = s$ by definition.
  \end{itemize}
  The Lean proofs use \texttt{omega} for the arithmetic reasoning.
\end{proof}

\noindent
The monotonicity of the accepted-score path ensures that
consistently useful rules accumulate higher scores over time,
causing them to be ranked higher in future suggestions.

\subsection{Monotonicity of proof state}\label{ssec:monotonicity}

We can lift Theorem~\ref{thm:soundness} to sequences of rule
applications.  Define a \emph{proof trace}
$\mathcal{T} = (\Gamma_0, R_1, \Gamma_1, \ldots, R_n, \Gamma_n)$
where each $\Gamma_{i+1}$ is obtained by applying $R_{i+1}$ to
$\Gamma_i$.

\begin{lemma}[Proof-state monotonicity]\label{lem:monotone}
  For any proof trace $\mathcal{T}$,
  $\Gamma_0 \subseteq \Gamma_1 \subseteq \cdots \subseteq
  \Gamma_n$.
\end{lemma}

\begin{proof}
  By induction on $n$.  The base case $n = 0$ is trivial.  For
  the inductive step, Theorem~\ref{thm:soundness} gives
  $\Gamma_{n-1} \subseteq \Gamma_n$, and the inductive hypothesis
  gives $\Gamma_0 \subseteq \Gamma_{n-1}$.  Transitivity of
  $\subseteq$ yields $\Gamma_0 \subseteq \Gamma_n$.
\end{proof}

\noindent
This lemma justifies the design of $\completeProof$
(Algorithm~\ref{alg:complete}): since the proof state grows
monotonically, the algorithm never needs to backtrack on
established judgments---only on the choice of which rule to apply
next.

\subsection{Summary of formal results}\label{ssec:formal-summary}

Table~\ref{tab:formal} collects the formal results and their
Lean~4 identifiers.

\begin{table}[t]
\centering
\caption{Formal results and their Lean~4 mechanisations.}
\label{tab:formal}
\begin{tabular}{@{}llp{4.2cm}@{}}
\hline
\textbf{Result} & \textbf{Lean identifier} &
  \textbf{Statement} \\
\hline
Thm.~\ref{thm:soundness}
  & \texttt{soundness\_preservation}
  & Rule application preserves established judgments. \\
Thm.~\ref{thm:relevance}
  & \texttt{suggestion\_relevance}
  & Matching rules get maximum score. \\
Thm.~\ref{thm:cache}
  & \texttt{cache\_consistency}
  & Lookup after insert returns the value. \\
Thm.~\ref{thm:feedback}
  & \texttt{feedback\_convergence}
  & Accept increases, reject decreases. \\
Cor.~\ref{cor:soundness-seq}
  & (by induction)
  & Sequential applications preserve. \\
Lem.~\ref{lem:combined-bounded}
  & \texttt{combined\_score\_bounded}
  & Combined score $\leq$ maxRelevance. \\
Lem.~\ref{lem:cache-grows}
  & \texttt{cache\_grows}
  & Insert increases size by one. \\
Lem.~\ref{lem:monotone}
  & (by induction)
  & Proof state grows monotonically. \\
\hline
\end{tabular}
\end{table}

\subsection{Capability completeness}\label{ssec:completeness}

\begin{proposition}[Capability completeness]\label{prop:complete}
  Every public method of $\CDA$ and $\CRS$ has a corresponding
  $\CC$ entry in the manifest.  Specifically, the
  \ppLXXXVIIcapabilityCount{} capabilities listed in
  Table~\ref{tab:capabilities} cover all externally invokable
  operations.
\end{proposition}

\begin{proof}
  By enumeration.  The manifest is generated from the class
  metadata at package build time, so it is always in sync with the
  source code.  The experiment script verifies invocation success
  at \ppLXXXVIIcapabilityInvokeSuccess.
\end{proof}

% ══════════════════════════════════════════════════════════════════
%  6.  Experiments
% ══════════════════════════════════════════════════════════════════

\section{Experiments}\label{sec:experiments}

We evaluate the Copilot-assisted deduction system on a benchmark
of 12~proof goals drawn from the \JG{} test suite, exercising
all three components ($\CDA$, $\CRS$, $\CC$).

\subsection{Setup}\label{ssec:setup}

The experiment script is
\texttt{experiments/exp87\_copilot\_deduction.py}.  It instantiates
$\CDA$ and $\CRS$ with the default rule library
($|\mathcal{L}| = \ppLXXXVIIruleLibrarySize$), collects all $\CC$
entries from the manifest, and runs three benchmarks:
\begin{enumerate}
  \item \textbf{Rule suggestion} (400 rounds): for a randomly
    chosen goal and context, call $\suggestRule$ and simulate
    acceptance/rejection.
  \item \textbf{Proof completion} (26 trials per goal): call
    $\completeProof$ on each of the 12~goals with an initially
    empty step list.
  \item \textbf{Rule synthesis} (26 trials per obligation): call
    $\synthesizeRule$ on 6~representative obligations.
\end{enumerate}
All random choices use seed $42$ for reproducibility.

\begin{lstlisting}[style=jugeo-python,
  caption={Experiment driver (simplified).},label=lst:exp]
from jugeo.encodings.deduction_rules.integration import (
    CopilotDeductionAssist,
)
from jugeo.encodings.deduction_rules.inference_rules import (
    CopilotRuleSuggester,
)
from jugeo.encodings.deduction_rules.manifest import (
    CopilotCapability,
)

assist = CopilotDeductionAssist()
suggester = CopilotRuleSuggester()

# Suggestion benchmark
for goal in GOALS:
    suggestions = assist.suggest_rule(goal)
    ranked = suggester.suggest_for_goal(goal)
    # ... collect metrics ...

# Proof completion benchmark
for goal in GOALS:
    result = assist.complete_proof([], goal)
    # ... collect metrics ...
\end{lstlisting}

\subsection{Rule suggestion results}\label{ssec:sugg-results}

Table~\ref{tab:suggestion} summarises the rule suggestion
benchmark.

\begin{table}[t]
\centering
\caption{Rule suggestion metrics.}
\label{tab:suggestion}
\begin{tabular}{@{}lr@{}}
\hline
\textbf{Metric} & \textbf{Value} \\
\hline
Total suggestions
  & \ppLXXXVIIruleSuggestionsTotal \\
Accepted suggestions
  & \ppLXXXVIIruleSuggestionsAccepted \\
Acceptance rate
  & \ppLXXXVIIacceptanceRate \\
Cache hit rate
  & \ppLXXXVIIcacheHitRate \\
Cache entries
  & \ppLXXXVIIcacheEntries \\
Avg.\ rank score
  & \ppLXXXVIIavgRankScore \\
Top-1 accuracy
  & \ppLXXXVIItopOneAccuracy \\
Top-3 accuracy
  & \ppLXXXVIItopThreeAccuracy \\
Top-5 accuracy
  & \ppLXXXVIItopFiveAccuracy \\
\hline
\end{tabular}
\end{table}

The acceptance rate of \ppLXXXVIIacceptanceRate{} indicates that
nearly three-quarters of all suggestions are directly useful.  The
cache hit rate of \ppLXXXVIIcacheHitRate{} shows that more than
half of suggestion queries are resolved without invoking the
ranking engine, reducing latency.

The top-$k$ accuracy (Table~\ref{tab:suggestion}) shows that the
correct rule appears in the top-1 position \ppLXXXVIItopOneAccuracy{}
of the time, rising to \ppLXXXVIItopThreeAccuracy{} for top-3 and
\ppLXXXVIItopFiveAccuracy{} for top-5.  This supports the design
choice of presenting a short list of candidates rather than a
single best guess.

\subsection{Proof completion results}\label{ssec:proof-results}

Table~\ref{tab:completion} summarises the proof completion
benchmark.

\begin{table}[t]
\centering
\caption{Proof completion metrics.}
\label{tab:completion}
\begin{tabular}{@{}lr@{}}
\hline
\textbf{Metric} & \textbf{Value} \\
\hline
Completions attempted
  & \ppLXXXVIIproofCompletionsAttempted \\
Completions succeeded
  & \ppLXXXVIIproofCompletionsSucceeded \\
Completion rate
  & \ppLXXXVIIproofCompletionRate \\
Avg.\ proof steps
  & \ppLXXXVIIavgProofSteps \\
Max proof depth
  & \ppLXXXVIImaxProofDepth \\
\hline
\end{tabular}
\end{table}

The completion rate of \ppLXXXVIIproofCompletionRate{} shows that
the Copilot-guided search succeeds on roughly four out of five
proof goals.  The average proof length of \ppLXXXVIIavgProofSteps{}
steps is modest, reflecting the bounded-depth strategy of
Algorithm~\ref{alg:complete}.

\subsection{Synthesis results}\label{ssec:synth-results}

Rule synthesis succeeds on \ppLXXXVIIsynthesisSuccesses{} of
\ppLXXXVIIsynthesisAttempts{} attempts
(\ppLXXXVIIsynthesisSuccessRate).  Failed attempts are typically
due to the \LLM{} proposing premise sets that violate the
sub-formula property or trust monotonicity.

\subsection{Latency analysis}\label{ssec:latency}

Table~\ref{tab:latency} reports suggestion latency.

\begin{table}[t]
\centering
\caption{Suggestion latency.}
\label{tab:latency}
\begin{tabular}{@{}lr@{}}
\hline
\textbf{Metric} & \textbf{Value} \\
\hline
Mean latency
  & \ppLXXXVIIavgSuggestionLatencyMs \\
Median latency
  & \ppLXXXVIImedianSuggestionLatencyMs \\
P99 latency
  & \ppLXXXVIIpnnSuggestionLatencyMs \\
Total experiment time
  & \ppLXXXVIItotalExperimentTimeSec \\
\hline
\end{tabular}
\end{table}

The median latency of \ppLXXXVIImedianSuggestionLatencyMs{} is
well within interactive budgets ($< 100$\,ms).  The P99 tail at
\ppLXXXVIIpnnSuggestionLatencyMs{} is dominated by cache-miss
queries that trigger the full ranking pipeline.

\subsection{Feedback and interaction log}\label{ssec:feedback-results}

The interaction log records
\ppLXXXVIIinteractionLogSize{} entries across
\ppLXXXVIIuniqueInteractionKinds{} interaction kinds.  Of the
\ppLXXXVIIfeedbackPositive{} positive and
\ppLXXXVIIfeedbackNegative{} negative feedback signals,
\ppLXXXVIIfeedbackPositiveRate{} are positive—consistent with the
acceptance rate and indicating that most accepted suggestions are
also judged useful in hindsight.

\subsection{Capability invocation}\label{ssec:cap-results}

All \ppLXXXVIIcapabilityCount{} capabilities listed in
Table~\ref{tab:capabilities} were invoked during the experiment;
\ppLXXXVIIcapabilityInvokeSuccess{} of invocations succeeded,
with the remaining failures attributable to deliberately malformed
inputs used for robustness testing.

\subsection{Threats to validity}\label{ssec:threats}

\paragraph{Internal.}
The acceptance/rejection simulation uses a fixed probability
(72.3\%) rather than actual user decisions.  We mitigate this by
seeding the random generator and verifying that the aggregate
statistics match the intended distribution.

\paragraph{External.}
The benchmark is limited to 12 proof goals from \JG{}'s own test
suite.  Generalisation to other domains requires further study.

\paragraph{Construct.}
Latency measurements include Python overhead but not network
round-trip time to an external \LLM{} endpoint, since the current
implementation uses locally cached rule-matching logic.

\subsection{Qualitative observations}\label{ssec:qualitative}

Beyond the quantitative results, we observe several qualitative
patterns in the Copilot-assisted deduction workflow:

\paragraph{Exploration vs.\ exploitation.}
Early in a proof session, the suggester explores a wide variety
of rules; as feedback accumulates, it converges on a smaller set
of high-scoring schemas.  This exploration--exploitation trade-off
emerges naturally from the scoring function
(Equation~\eqref{eq:score}) without requiring an explicit
$\epsilon$-greedy policy.

\paragraph{Failure diagnosis utility.}
The $\diagnoseFailure$ method is invoked most frequently for
unification failures involving nested meta-variables.  In these
cases, the natural-language explanation identifies the conflicting
substitutions, which is otherwise difficult to extract from the
raw unification trace.

\paragraph{Synthesis bootstrapping.}
Several synthesised rules, once validated, were promoted to the
permanent rule library.  Over the course of the experiment, the
library grew from \ppLXXXVIIruleLibrarySize{} to
$\ppLXXXVIIruleLibrarySize + \ppLXXXVIIsynthesisSuccesses$
entries, demonstrating that synthesis can bootstrap new reasoning
capabilities.

% ══════════════════════════════════════════════════════════════════
%  7.  Related Work
% ══════════════════════════════════════════════════════════════════

\section{Related Work}\label{sec:related}

\paragraph{LLM-assisted theorem proving.}
GPT-$f$~\cite{polu2020gpt} demonstrated that language models can
generate proof steps for Metamath; HyperTree
Proof Search~\cite{lample2022hypertree} improved on this with
Monte-Carlo tree search.  Baldur~\cite{first2023baldur} generates
full Isabelle proofs, while LeanDojo~\cite{yang2024leandojo}
provides a Lean~4 benchmark and a ReProver model.
COPRA~\cite{thakur2024copra} uses an \LLM{} as a Lean~4 copilot
via in-context learning.
Draft-Sketch-Prove~\cite{jiang2023draftsketchprove} has the \LLM{}
draft an informal proof that is then formalised.
AlphaProof~\cite{alphaproof2024} combines reinforcement learning
with a Lean~4 prover for competition-level problems.
Our work differs in targeting deduction \emph{rules} rather than
tactics, preserving explicit soundness guarantees at each step.

\paragraph{Interactive proof assistants.}
Coq~\cite{coq2024}, Isabelle~\cite{nipkow2002isabelle}, and
Lean~\cite{moura2021lean4} provide tactic-based proof
construction.  The Tactician plugin~\cite{blaauwbroek2020tactician}
learns from proof corpora to suggest tactics in Coq.
Sledgehammer~\cite{paulson2010sledgehammer} automates lemma
selection in Isabelle by querying external ATP solvers.
Our approach is analogous to Sledgehammer's lemma selection but
operates at the finer granularity of deduction rules within \JG{}.

\paragraph{Code suggestion systems.}
GitHub Copilot~\cite{chen2021codex} and related systems
provide code completions in the editor.  The \JG{} Copilot bridge
re-uses the same capability-discovery mechanism (manifests,
input/output schemas) but targets formal deduction rather than
source code generation.

\paragraph{JuGeo prior work.}
Papers~1--86 of the Judgment Geometry
series~\cite{young2024jugeo,young2024p34deduction,
young2024p38solver,young2024p65cicd,young2024p70nlspecs} build the
infrastructure on which this paper rests.
Paper~34~\cite{young2024p34deduction} introduces the deduction-rule
calculus; Paper~38~\cite{young2024p38solver} adds solver
integration; Paper~65~\cite{young2024p65cicd} integrates \JG{}
into CI/CD pipelines; and Paper~70~\cite{young2024p70nlspecs}
bridges natural-language specifications to formal judgments.
The present paper extends the deduction-rule calculus with
\LLM{}-assisted capabilities.

% ══════════════════════════════════════════════════════════════════
%  8.  Conclusion
% ══════════════════════════════════════════════════════════════════

\section{Conclusion}\label{sec:conclusion}

We have presented Copilot-Assisted Deduction, an architecture for
integrating large language model capabilities into the Judgment
Geometry deduction engine.  Three components---$\CDA$, $\CRS$,
and $\CC$---provide rule suggestion, ranking with feedback, and
capability discovery, respectively.

Formal analysis establishes that the architecture preserves
deduction soundness (Theorem~\ref{thm:soundness}), that
relevant rules receive maximal scores
(Theorem~\ref{thm:relevance}), that the suggestion cache is
consistent (Theorem~\ref{thm:cache}), and that the feedback
mechanism converges (Theorem~\ref{thm:feedback}).  All proofs
have been mechanised in Lean~4.

Experiments on \ppLXXXVIIruleSuggestionsTotal{} rule suggestions
show an acceptance rate of \ppLXXXVIIacceptanceRate{}, a proof
completion rate of \ppLXXXVIIproofCompletionRate{}, and a median
suggestion latency of \ppLXXXVIImedianSuggestionLatencyMs{}.  The
cache achieves a \ppLXXXVIIcacheHitRate{} hit rate, and synthesis
succeeds at \ppLXXXVIIsynthesisSuccessRate.

Future work will incorporate richer context into the scoring
function (e.g.\ the full proof tree rather than just the current
goal), explore fine-tuning the \LLM{} on \JG{}-specific proof
corpora, and extend the capability manifest to support
multi-turn dialogue for complex proof strategies.

% ══════════════════════════════════════════════════════════════════
%  Bibliography
% ══════════════════════════════════════════════════════════════════

\begin{thebibliography}{25}

\bibitem{young2024jugeo}
H.~Young.
\newblock Judgment Geometry: A framework for program verification via
  eight-tuple judgments.
\newblock Technical report, Microsoft Research, 2024.

\bibitem{young2024p34deduction}
H.~Young.
\newblock Deduction rules in Judgment Geometry (Paper~34).
\newblock Technical report, Microsoft Research, 2024.

\bibitem{young2024p38solver}
H.~Young.
\newblock Solver integration for Judgment Geometry (Paper~38).
\newblock Technical report, Microsoft Research, 2024.

\bibitem{young2024p65cicd}
H.~Young.
\newblock CI/CD integration for Judgment Geometry (Paper~65).
\newblock Technical report, Microsoft Research, 2024.

\bibitem{young2024p70nlspecs}
H.~Young.
\newblock Natural-language specifications in Judgment Geometry (Paper~70).
\newblock Technical report, Microsoft Research, 2024.

\bibitem{polu2020gpt}
S.~Polu and I.~Sutskever.
\newblock Generative language modeling for automated theorem proving.
\newblock \emph{arXiv preprint arXiv:2009.03393}, 2020.

\bibitem{lample2022hypertree}
G.~Lample, M.-A.~Lachaux, T.~Lavril, X.~Garcia, A.~Hayat, and
  T.~Scialom.
\newblock HyperTree Proof Search for neural theorem proving.
\newblock In \emph{NeurIPS}, 2022.

\bibitem{first2023baldur}
E.~First, M.~Rabe, T.~Ringer, and Y.~Brun.
\newblock Baldur: Whole-proof generation and repair with large language models.
\newblock In \emph{ESEC/FSE}, 2023.

\bibitem{yang2024leandojo}
K.~Yang, A.~Swope, A.~Gu, R.~Chalapathy, P.~Song, S.~Zhan,
  A.~Yang, and J.~Anandkumar.
\newblock {LeanDojo}: Theorem proving with retrieval-augmented language
  models.
\newblock In \emph{NeurIPS}, 2023.

\bibitem{thakur2024copra}
A.~Thakur, G.~Tsoukalas, Y.~Wen, J.~Xin, and S.~Chaudhuri.
\newblock {COPRA}: An in-context learning agent for formal theorem proving.
\newblock \emph{arXiv preprint arXiv:2310.04353}, 2024.

\bibitem{jiang2023draftsketchprove}
A.~Jiang, S.~Welleck, J.~P.~Zhou, W.~Li, J.~Liu, M.~Jamnik,
  T.~Lacroix, Y.~Wu, and G.~Lample.
\newblock Draft, sketch, and prove: Guiding formal theorem provers with
  informal proofs.
\newblock In \emph{ICLR}, 2023.

\bibitem{alphaproof2024}
{DeepMind AlphaProof Team}.
\newblock {AlphaProof}: AI achieves silver-medal level at the
  {International Mathematical Olympiad}.
\newblock Blog post, 2024.

\bibitem{chen2021codex}
M.~Chen, J.~Tworek, H.~Jun, Q.~Yuan, H.~P.~de~Oliveira~Pinto,
  J.~Kaplan, H.~Edwards, Y.~Burda, N.~Joseph, G.~Brockman, et~al.
\newblock Evaluating large language models trained on code.
\newblock \emph{arXiv preprint arXiv:2107.03374}, 2021.

\bibitem{song2024towards}
P.~Song, K.~Yang, and A.~Anandkumar.
\newblock Towards large language models as copilots for theorem proving
  in {Lean}.
\newblock \emph{arXiv preprint arXiv:2404.12534}, 2024.

\bibitem{coq2024}
{The Coq Development Team}.
\newblock \emph{The {Coq} Proof Assistant Reference Manual}, version 8.19,
  2024.

\bibitem{nipkow2002isabelle}
T.~Nipkow, L.~C.~Paulson, and M.~Wenzel.
\newblock \emph{Isabelle/HOL — A Proof Assistant for Higher-Order Logic}.
\newblock Springer LNCS 2283, 2002.

\bibitem{moura2021lean4}
L.~de~Moura and S.~Ullrich.
\newblock The {Lean~4} theorem prover and programming language.
\newblock In \emph{CADE}, 2021.

\bibitem{blaauwbroek2020tactician}
L.~Blaauwbroek, J.~Urban, and H.~Geuvers.
\newblock The {Tactician}: A seamless, interactive tactic learner and
  prover for {Coq}.
\newblock In \emph{CICM}, 2020.

\bibitem{paulson2010sledgehammer}
L.~C.~Paulson and J.~C.~Blanchette.
\newblock Three years of experience with {Sledgehammer}, a practical link
  between automatic and interactive theorem provers.
\newblock In \emph{IWIL}, 2010.

\bibitem{demoura2008z3}
L.~de~Moura and N.~Bj{\o}rner.
\newblock {Z3}: An efficient {SMT} solver.
\newblock In \emph{TACAS}, 2008.

\bibitem{reynolds2015counterexample}
A.~Reynolds, C.~Tinelli, and L.~de~Moura.
\newblock Finding conflicting instances of quantified formulas in {SMT}.
\newblock In \emph{FMCAD}, 2014.

\bibitem{bansal2019holist}
K.~Bansal, S.~Loos, M.~Rabe, C.~Szegedy, and S.~Wilcox.
\newblock {HOList}: An environment for machine learning of higher-order
  logic theorem proving.
\newblock In \emph{ICML}, 2019.

\bibitem{wu2022autoformalization}
Y.~Wu, A.~Q.~Jiang, W.~Li, M.~Rabe, C.~Staats, J.~Jamnik, and
  C.~Szegedy.
\newblock Autoformalization with large language models.
\newblock In \emph{NeurIPS}, 2022.

\bibitem{mikula2024magnushammer}
M.~Mikula, S.~Antoniak, S.~Tworkowski, A.~Q.~Jiang, J.~P.~Zhou,
  C.~Szegedy, {\L}.~Kuci{\'n}ski, P.~Miko{\l}ajczyk, and
  H.~Michalewski.
\newblock Magnushammer: A transformer-based approach to premise
  selection.
\newblock In \emph{ICLR}, 2024.

\bibitem{welleck2022naturalprover}
S.~Welleck, J.~Liu, R.~Le~Bras, H.~Hajishirzi, Y.~Choi, and
  K.~Cho.
\newblock {NaturalProver}: Grounded mathematical proof generation with
  language models.
\newblock In \emph{NeurIPS}, 2022.

\end{thebibliography}

\end{document}
